\documentclass[11pt]{article}
\newcommand{\DocNumber}{VOLUME 02}
\newcommand{\DocTitle}{C++23 Standard Library Reference}
\newcommand{\DocSubtitle}{Facilities and semantic contracts, clauses 16--33}
\newcommand{\DocSummary}{A full library-clause map with header-aware examples, API selection guidance and review questions.}
% Shared, pdfLaTeX-only design system for the C++23 documentation suite.
\usepackage[T1]{fontenc}
\usepackage[utf8]{inputenc}
\usepackage{lmodern}
\usepackage[english]{babel}
\usepackage[letterpaper,margin=0.82in,headheight=16pt,headsep=17pt]{geometry}
\usepackage{microtype}
\usepackage{amsmath,amssymb}
\usepackage{array,booktabs,tabularx,longtable}
\usepackage[table]{xcolor}
\usepackage[most]{tcolorbox}
\usepackage{listings}
\usepackage{fancyhdr}
\usepackage{titlesec}
\usepackage{enumitem}
\usepackage{hyperref}
\usepackage{needspace}
\usepackage{etoolbox}
\definecolor{ink}{HTML}{13253B}
\definecolor{teal}{HTML}{006E7F}
\definecolor{muted}{HTML}{5A6C7C}
\definecolor{fog}{HTML}{EDF4F6}
\definecolor{sand}{HTML}{FAF3E8}
\definecolor{rule}{HTML}{C9D7DF}
\definecolor{codebg}{HTML}{F7F9FB}
\definecolor{codecomment}{HTML}{47715D}
\hypersetup{colorlinks=true,linkcolor=teal,urlcolor=teal,pdfauthor={Independent C++23 engineering documentation},pdftitle={C++23 Technical Suite}}
\setlength{\parindent}{0pt}
\setlength{\parskip}{6pt}
\setcounter{tocdepth}{2}
\renewcommand{\arraystretch}{1.18}
\titleformat{\section}{\Large\bfseries\color{ink}}{\thesection}{0.7em}{}
\titleformat{\subsection}{\large\bfseries\color{teal}}{\thesubsection}{0.55em}{}
\titlespacing*{\section}{0pt}{17pt}{7pt}
\titlespacing*{\subsection}{0pt}{11pt}{4pt}
\pagestyle{fancy}
\fancyhf{}
\fancyhead[L]{\footnotesize\bfseries\color{teal}C++23 / ISO 14882:2024}
\fancyhead[R]{\footnotesize\color{muted}\DocNumber}
\fancyfoot[L]{\scriptsize\color{muted}Independent engineering companion, not an ISO publication}
\fancyfoot[R]{\scriptsize\color{muted}\thepage}
\renewcommand{\headrulewidth}{0.35pt}
\lstdefinestyle{cppstyle}{language=C++,basicstyle=\ttfamily\footnotesize,keywordstyle=\color{teal}\bfseries,commentstyle=\color{codecomment},stringstyle=\color{ink},backgroundcolor=\color{codebg},frame=single,rulecolor=\color{rule},breaklines=true,breakatwhitespace=true,columns=fullflexible,keepspaces=true,showstringspaces=false,tabsize=2,numbers=none,upquote=true,xleftmargin=3pt,xrightmargin=3pt,aboveskip=6pt,belowskip=8pt}
\lstset{style=cppstyle}
\newtcolorbox{standardbox}{enhanced,breakable,colback=fog,colframe=teal!35!white,boxrule=0.55pt,arc=1.4mm,left=3mm,right=3mm,top=2mm,bottom=2mm,before skip=7pt,after skip=8pt}
\newtcolorbox{practicebox}{enhanced,breakable,colback=sand,colframe=ink!20!white,boxrule=0.55pt,arc=1.4mm,left=3mm,right=3mm,top=2mm,bottom=2mm,before skip=7pt,after skip=8pt}
\newcommand{\anchorref}[1]{\par\smallskip{\footnotesize\textbf{Standard locator:} \texttt{#1}}\par}
\newcommand{\standard}[1]{\begin{standardbox}\textbf{Standard-grounded interpretation.} #1\end{standardbox}}
\newcommand{\practice}[1]{\begin{practicebox}\textbf{Engineering recommendation (not a normative ISO rule).} #1\end{practicebox}}
\newcommand{\takeaway}[1]{\par\textbf{Review question.} #1\par}
\newcommand{\docfront}{%
\begin{titlepage}
\vspace*{25mm}
{\color{teal}\bfseries\large ENGINEERING REFERENCE SERIES\par}
\vspace{12mm}
{\raggedright\color{ink}\bfseries\fontsize{24}{30}\selectfont\DocTitle\par}
\vspace{7mm}
{\color{muted}\fontsize{15}{19}\selectfont\DocSubtitle\par}
\vspace{12mm}
\noindent\color{teal}\rule{\linewidth}{2pt}\par
\vspace{9mm}
{\large\DocSummary\par}
\vfill
\begin{tabularx}{\linewidth}{@{}lX@{}}
\textbf{Volume} & \DocNumber\\
\textbf{Standards basis} & ISO/IEC 14882:2024 (seventh edition, C++23)\\
\textbf{Document type} & Original interpretive and instructional material\\
\textbf{Date} & 16 October 2026\\
\textbf{Toolchain} & pdfLaTeX; illustrative C++23 examples require separate compiler support\\
\end{tabularx}
\vspace{8mm}
{\footnotesize This document is not an authorized copy of, replacement for, or amendment to ISO/IEC 14882:2024. The original standard controls in the event of disagreement.}\par
\end{titlepage}
\pagenumbering{roman}
\tableofcontents
\clearpage
\pagenumbering{arabic}
\section*{How to use this volume}\addcontentsline{toc}{section}{How to use this volume}
The reference is \emph{ISO/IEC 14882:2024, Programming languages --- C++}. Section locators and square-bracketed stable names identify where to verify precise wording; they are not quotations. \textbf{Standard-grounded interpretation} boxes summarize language or library semantics. \textbf{Engineering recommendation} boxes contain optional design, security, portability, or process guidance. Code is original and illustrative; individual library facilities may not yet be implemented by every vendor. This suite does not grant rights to reproduce or distribute the licensed standard.\par\medskip
}
\newcommand{\docend}{\clearpage\section*{Reference and traceability}\addcontentsline{toc}{section}{Reference and traceability}
\textbf{Primary source.} ISO/IEC 14882:2024(en), \emph{Programming languages --- C++}, seventh edition, October 2024. Use its numbered clauses, stable names, Annex C (informative compatibility), and Annex D (compatibility features) for authoritative details.\par
\textbf{Limits.} Clause-level summaries and exercises in this document are selective explanations, not exhaustive reproductions of all semantic constraints, library mandates, complexity clauses, and overload participation requirements. A successful compiler test demonstrates behavior of that specific implementation and configuration, not universal conformance.\par
\textbf{Use of contemporary engineering tools.} Compilers, sanitizers, analyzers, CI systems, repository policies, and vulnerability taxonomies are implementation guidance rather than language-standard requirements.\par}

\begin{document}
\docfront
\section{Library contract interpretation}
Library clauses define requirements across operations: preconditions, effects, return values, throws, remarks and complexity. A container API call can compile and still violate its semantic precondition. Feature availability is an implementation property, not implied merely by selecting -std=c++23.

\small
\begin{longtable}{@{}p{.25\linewidth}p{.69\linewidth}@{}}
\toprule
\textbf{Question} & \textbf{Review method}\\
\midrule
\endfirsthead
\toprule
\textbf{Question} & \textbf{Review method}\\
\midrule
\endhead
Is it in the 2024 library? & Locate the clause/header in 16--33.\\
\addlinespace[2pt]
Is the implementation complete? & Check the vendor status and the corresponding feature-test macro.\\
\addlinespace[2pt]
Can this call fail? & Inspect preconditions, exception clauses, and returned error values.\\
\addlinespace[2pt]
Can a reference dangle? & Evaluate ownership and iterator/view invalidation.\\
\addlinespace[2pt]
Is an optimization legal? & Inspect documented complexity and observable behavior guarantees.\\
\addlinespace[2pt]
\bottomrule
\end{longtable}
\normalsize
\section{Library model and requirements}
\anchorref{16}
Core scope: The library specification defines required expressions, semantic preconditions, exception effects, complexity, and implementer latitude. Library contracts are not the same as style conventions.

\standard{ISO/IEC 14882:2024 defines the referenced facility in Clause 16. Precise overload participation, mandates and complexity require consulting the cited clauses.}
\begin{itemize}[leftmargin=*,itemsep=1pt,topsep=3pt]
\item Identify the facility header and type constraints.
\item Check preconditions and iterator/reference invalidation where applicable.
\item Test representative and boundary cases on the selected implementation.
\end{itemize}
\begin{lstlisting}[language=C++]
#include <version>
#ifdef __cpp_lib_expected
// Inspect feature macro value before using <expected>.
#endif
\end{lstlisting}
\practice{Confirm header availability, feature macros and semantic requirements before selecting an API.}
\takeaway{What invariant must hold before invoking this facility?}
\section{Language support utilities}
\anchorref{17}
Core scope: This clause supplies fundamentals such as type information, source location, dynamic allocation interfaces, coroutine support and special runtime types.

\standard{ISO/IEC 14882:2024 defines the referenced facility in Clause 17. Precise overload participation, mandates and complexity require consulting the cited clauses.}
\begin{itemize}[leftmargin=*,itemsep=1pt,topsep=3pt]
\item Identify the facility header and type constraints.
\item Check preconditions and iterator/reference invalidation where applicable.
\item Test representative and boundary cases on the selected implementation.
\end{itemize}
\begin{lstlisting}[language=C++]
#include <source_location>
#include <string_view>
void log(std::string_view message,
         std::source_location from = std::source_location::current());
\end{lstlisting}
\practice{Differentiate source\_location capture from authorization or trusted audit metadata.}
\takeaway{What invariant must hold before invoking this facility?}
\section{Concepts library}
\anchorref{18}
Core scope: Library concepts such as integral, regular and invocable express syntactic and semantic requirements used by generic facilities.

\standard{ISO/IEC 14882:2024 defines the referenced facility in Clause 18. Precise overload participation, mandates and complexity require consulting the cited clauses.}
\begin{itemize}[leftmargin=*,itemsep=1pt,topsep=3pt]
\item Identify the facility header and type constraints.
\item Check preconditions and iterator/reference invalidation where applicable.
\item Test representative and boundary cases on the selected implementation.
\end{itemize}
\begin{lstlisting}[language=C++]
#include <concepts>
template<std::regular T> struct Store { T item; };
\end{lstlisting}
\practice{A satisfied constraint can still fail semantic modeling requirements such as equality preservation.}
\takeaway{What invariant must hold before invoking this facility?}
\section{Diagnostics library}
\anchorref{19}
Core scope: Exception and error types, assertion support, system errors and stacktraces provide standardized diagnostic vocabulary.

\standard{ISO/IEC 14882:2024 defines the referenced facility in Clause 19. Precise overload participation, mandates and complexity require consulting the cited clauses.}
\begin{itemize}[leftmargin=*,itemsep=1pt,topsep=3pt]
\item Identify the facility header and type constraints.
\item Check preconditions and iterator/reference invalidation where applicable.
\item Test representative and boundary cases on the selected implementation.
\end{itemize}
\begin{lstlisting}[language=C++]
#include <system_error>
#include <string>
std::error_code ec;
std::string msg = ec.message();
\end{lstlisting}
\practice{Use error\_code for recoverable platform failures, and avoid asserting on untrusted input.}
\takeaway{What invariant must hold before invoking this facility?}
\section{Memory management library}
\anchorref{20}
Core scope: Allocators, unique/shared/weak pointers, ownership utilities and memory resources implement storage and ownership abstractions.

\standard{ISO/IEC 14882:2024 defines the referenced facility in Clause 20. Precise overload participation, mandates and complexity require consulting the cited clauses.}
\begin{itemize}[leftmargin=*,itemsep=1pt,topsep=3pt]
\item Identify the facility header and type constraints.
\item Check preconditions and iterator/reference invalidation where applicable.
\item Test representative and boundary cases on the selected implementation.
\end{itemize}
\begin{lstlisting}[language=C++]
#include <memory>
auto owned = std::make_unique<int>(42);
int answer = *owned;
\end{lstlisting}
\practice{Prefer unique ownership by default; audit shared ownership cycles and lifetime of polymorphic memory resources.}
\takeaway{What invariant must hold before invoking this facility?}
\section{Metaprogramming library}
\anchorref{21}
Core scope: Type traits and integer sequences allow compile-time transformation and inspection of properties.

\standard{ISO/IEC 14882:2024 defines the referenced facility in Clause 21. Precise overload participation, mandates and complexity require consulting the cited clauses.}
\begin{itemize}[leftmargin=*,itemsep=1pt,topsep=3pt]
\item Identify the facility header and type constraints.
\item Check preconditions and iterator/reference invalidation where applicable.
\item Test representative and boundary cases on the selected implementation.
\end{itemize}
\begin{lstlisting}[language=C++]
#include <type_traits>
static_assert(std::is_integral_v<unsigned long>);
static_assert(!std::is_reference_v<int>);
\end{lstlisting}
\practice{Prefer standard type traits over home-grown trait specializations; document type-trait preconditions.}
\takeaway{What invariant must hold before invoking this facility?}
\section{General utilities library}
\anchorref{22}
Core scope: optional, variant, expected, tuples, function wrappers, formatting, bit operations and execution policies support value-oriented programming.

\standard{ISO/IEC 14882:2024 defines the referenced facility in Clause 22. Precise overload participation, mandates and complexity require consulting the cited clauses.}
\begin{itemize}[leftmargin=*,itemsep=1pt,topsep=3pt]
\item Identify the facility header and type constraints.
\item Check preconditions and iterator/reference invalidation where applicable.
\item Test representative and boundary cases on the selected implementation.
\end{itemize}
\begin{lstlisting}[language=C++]
#include <expected>
#include <string>
std::expected<int, std::string> parse(bool ok) {
    if (!ok) return std::unexpected(std::string{"invalid"});
    return 7;
}
\end{lstlisting}
\practice{Use optional for absence and expected for operations with a meaningful error alternative; never treat the two as synonyms.}
\takeaway{What invariant must hold before invoking this facility?}
\section{Strings library}
\anchorref{23}
Core scope: basic\_string owns characters; basic\_string\_view supplies a non-owning view. Termination and lifetime conditions differ.

\standard{ISO/IEC 14882:2024 defines the referenced facility in Clause 23. Precise overload participation, mandates and complexity require consulting the cited clauses.}
\begin{itemize}[leftmargin=*,itemsep=1pt,topsep=3pt]
\item Identify the facility header and type constraints.
\item Check preconditions and iterator/reference invalidation where applicable.
\item Test representative and boundary cases on the selected implementation.
\end{itemize}
\begin{lstlisting}[language=C++]
#include <string>
#include <string_view>
std::string_view prefix(std::string_view s) {
    return s.substr(0, s.find(':'));
}
\end{lstlisting}
\practice{Never return a view into a temporary string; prefer explicit length and encoding contracts.}
\takeaway{What invariant must hold before invoking this facility?}
\section{Containers library}
\anchorref{24}
Core scope: Sequence, associative, unordered, adaptor and view/container facilities offer different iterator stability and algorithmic bounds.

\standard{ISO/IEC 14882:2024 defines the referenced facility in Clause 24. Precise overload participation, mandates and complexity require consulting the cited clauses.}
\begin{itemize}[leftmargin=*,itemsep=1pt,topsep=3pt]
\item Identify the facility header and type constraints.
\item Check preconditions and iterator/reference invalidation where applicable.
\item Test representative and boundary cases on the selected implementation.
\end{itemize}
\begin{lstlisting}[language=C++]
#include <vector>
#include <algorithm>
std::vector<int> ids{3, 1, 2};
std::ranges::sort(ids);
\end{lstlisting}
\practice{Select by element lifetime, invalidation behavior, lookup profile and mutation pattern, not naming preference.}
\takeaway{What invariant must hold before invoking this facility?}
\section{Iterators library}
\anchorref{25}
Core scope: Iterator concepts, categories, adaptors and range access define traversal contracts for algorithms.

\standard{ISO/IEC 14882:2024 defines the referenced facility in Clause 25. Precise overload participation, mandates and complexity require consulting the cited clauses.}
\begin{itemize}[leftmargin=*,itemsep=1pt,topsep=3pt]
\item Identify the facility header and type constraints.
\item Check preconditions and iterator/reference invalidation where applicable.
\item Test representative and boundary cases on the selected implementation.
\end{itemize}
\begin{lstlisting}[language=C++]
#include <iterator>
#include <vector>
static_assert(std::random_access_iterator<std::vector<int>::iterator>);
\end{lstlisting}
\practice{Do not retain iterators across operations without checking invalidation requirements.}
\takeaway{What invariant must hold before invoking this facility?}
\section{Ranges library}
\anchorref{26}
Core scope: Ranges model sequences and adaptors, including borrowed-range distinctions and newer generators. Lazy views defer operations until iteration.

\standard{ISO/IEC 14882:2024 defines the referenced facility in Clause 26. Precise overload participation, mandates and complexity require consulting the cited clauses.}
\begin{itemize}[leftmargin=*,itemsep=1pt,topsep=3pt]
\item Identify the facility header and type constraints.
\item Check preconditions and iterator/reference invalidation where applicable.
\item Test representative and boundary cases on the selected implementation.
\end{itemize}
\begin{lstlisting}[language=C++]
#include <ranges>
#include <vector>
auto positive = std::vector<int>{-1, 3, 4}
    | std::views::filter([](int x){ return x > 0; });
\end{lstlisting}
\practice{Trace ownership through views, avoid dangling range factories, and test single-pass assumptions.}
\takeaway{What invariant must hold before invoking this facility?}
\section{Algorithms library}
\anchorref{27}
Core scope: Non-modifying, mutating, ordering, numeric and parallel algorithms specify preconditions and complexity guarantees.

\standard{ISO/IEC 14882:2024 defines the referenced facility in Clause 27. Precise overload participation, mandates and complexity require consulting the cited clauses.}
\begin{itemize}[leftmargin=*,itemsep=1pt,topsep=3pt]
\item Identify the facility header and type constraints.
\item Check preconditions and iterator/reference invalidation where applicable.
\item Test representative and boundary cases on the selected implementation.
\end{itemize}
\begin{lstlisting}[language=C++]
#include <algorithm>
#include <vector>
std::vector<int> v{5, 2, 9};
std::ranges::sort(v);
\end{lstlisting}
\practice{Check iterator validity and comparator strict weak ordering; separate logical determinism from parallel scheduling.}
\takeaway{What invariant must hold before invoking this facility?}
\section{Numerics library}
\anchorref{28}
Core scope: Random engines and distributions, numeric arrays and floating-point interfaces have well-defined but different contracts.

\standard{ISO/IEC 14882:2024 defines the referenced facility in Clause 28. Precise overload participation, mandates and complexity require consulting the cited clauses.}
\begin{itemize}[leftmargin=*,itemsep=1pt,topsep=3pt]
\item Identify the facility header and type constraints.
\item Check preconditions and iterator/reference invalidation where applicable.
\item Test representative and boundary cases on the selected implementation.
\end{itemize}
\begin{lstlisting}[language=C++]
#include <numeric>
#include <vector>
std::vector<int> n{1, 2, 3};
int total = std::accumulate(n.begin(), n.end(), 0);
\end{lstlisting}
\practice{Do not use standard pseudorandom engines as cryptographic randomness; bound numerical error separately.}
\takeaway{What invariant must hold before invoking this facility?}
\section{Time library}
\anchorref{29}
Core scope: chrono distinguishes clocks, durations, time points, civil dates and time zones.

\standard{ISO/IEC 14882:2024 defines the referenced facility in Clause 29. Precise overload participation, mandates and complexity require consulting the cited clauses.}
\begin{itemize}[leftmargin=*,itemsep=1pt,topsep=3pt]
\item Identify the facility header and type constraints.
\item Check preconditions and iterator/reference invalidation where applicable.
\item Test representative and boundary cases on the selected implementation.
\end{itemize}
\begin{lstlisting}[language=C++]
#include <chrono>
auto t0 = std::chrono::steady_clock::now();
auto t1 = std::chrono::steady_clock::now();
auto elapsed = t1 - t0;
\end{lstlisting}
\practice{Use steady\_clock for intervals, wall clocks for timestamps, and explicit zone handling for civil scheduling.}
\takeaway{What invariant must hold before invoking this facility?}
\section{Localization library}
\anchorref{30}
Core scope: Locale facets, collation, classification and conversions affect text and numeric presentation.

\standard{ISO/IEC 14882:2024 defines the referenced facility in Clause 30. Precise overload participation, mandates and complexity require consulting the cited clauses.}
\begin{itemize}[leftmargin=*,itemsep=1pt,topsep=3pt]
\item Identify the facility header and type constraints.
\item Check preconditions and iterator/reference invalidation where applicable.
\item Test representative and boundary cases on the selected implementation.
\end{itemize}
\begin{lstlisting}[language=C++]
#include <locale>
#include <sstream>
std::ostringstream stream;
stream.imbue(std::locale::classic());
\end{lstlisting}
\practice{Keep wire formats locale-independent; document locale when interpreting human-facing input.}
\takeaway{What invariant must hold before invoking this facility?}
\section{Input/output library}
\anchorref{31}
Core scope: Streams, formatting, filesystem, span streams and printing encode different buffering, error, and resource semantics.

\standard{ISO/IEC 14882:2024 defines the referenced facility in Clause 31. Precise overload participation, mandates and complexity require consulting the cited clauses.}
\begin{itemize}[leftmargin=*,itemsep=1pt,topsep=3pt]
\item Identify the facility header and type constraints.
\item Check preconditions and iterator/reference invalidation where applicable.
\item Test representative and boundary cases on the selected implementation.
\end{itemize}
\begin{lstlisting}[language=C++]
#include <sstream>
#include <string>
std::istringstream in{"42"};
int value{};
if (!(in >> value)) { /* handle parse failure */ }
\end{lstlisting}
\practice{Check stream state and filesystem failures; separate user-controlled paths from trusted roots.}
\takeaway{What invariant must hold before invoking this facility?}
\section{Regular expressions library}
\anchorref{32}
Core scope: regex facilities provide search, matching, replacement and a specified modified ECMAScript grammar.

\standard{ISO/IEC 14882:2024 defines the referenced facility in Clause 32. Precise overload participation, mandates and complexity require consulting the cited clauses.}
\begin{itemize}[leftmargin=*,itemsep=1pt,topsep=3pt]
\item Identify the facility header and type constraints.
\item Check preconditions and iterator/reference invalidation where applicable.
\item Test representative and boundary cases on the selected implementation.
\end{itemize}
\begin{lstlisting}[language=C++]
#include <regex>
#include <string>
const std::regex digits{R"([0-9]+)"};
bool ok = std::regex_match(std::string{"123"}, digits);
\end{lstlisting}
\practice{Treat regex input size and complexity as operational risk; test adversarial inputs and prefer bounded parsers when appropriate.}
\takeaway{What invariant must hold before invoking this facility?}
\section{Concurrency support library}
\anchorref{33}
Core scope: Threads, atomics, synchronization, futures and coordination types specify data-race and ordering rules.

\standard{ISO/IEC 14882:2024 defines the referenced facility in Clause 33. Precise overload participation, mandates and complexity require consulting the cited clauses.}
\begin{itemize}[leftmargin=*,itemsep=1pt,topsep=3pt]
\item Identify the facility header and type constraints.
\item Check preconditions and iterator/reference invalidation where applicable.
\item Test representative and boundary cases on the selected implementation.
\end{itemize}
\begin{lstlisting}[language=C++]
#include <atomic>
std::atomic<unsigned> count{0};
void tick() { count.fetch_add(1, std::memory_order_relaxed); }
\end{lstlisting}
\practice{Prefer scoped synchronization; write down happens-before relationships and cancellation ownership.}
\takeaway{What invariant must hold before invoking this facility?}
\section{Library facility selection matrix}
\small
\begin{longtable}{@{}p{.18\linewidth}p{.40\linewidth}p{.34\linewidth}@{}}
\toprule
\textbf{Need} & \textbf{Primary candidates} & \textbf{Primary hazard}\\
\midrule
\endfirsthead
\toprule
\textbf{Need} & \textbf{Primary candidates} & \textbf{Primary hazard}\\
\midrule
\endhead
Optional value & optional & Calling value without engagement\\
\addlinespace[2pt]
Value or structured error & expected & Ignoring error state\\
\addlinespace[2pt]
Non-owning character sequence & string\_view & Dangling storage\\
\addlinespace[2pt]
Non-owning contiguous elements & span & Owner destruction or reallocation\\
\addlinespace[2pt]
Single owner & unique\_ptr & Escaping raw aliases\\
\addlinespace[2pt]
Shared owner & shared\_ptr/weak\_ptr & Cycles and contention\\
\addlinespace[2pt]
Lazy range pipeline & views & Deferred access to invalid storage\\
\addlinespace[2pt]
Correlated parallel state & mutex/atomic & Data race or deadlock\\
\addlinespace[2pt]
Human-readable output & format/print & Portability and unsupported formatter\\
\addlinespace[2pt]
\bottomrule
\end{longtable}
\normalsize
\docend
\end{document}
